\chapter{Paronychia}

\section*{Definition}
Paronychia is an inflammation of the nail folds (paronychium), presenting as painful erythema and swelling of the proximal or lateral nail fold. It is classified as acute (<6 weeks, usually infectious) or chronic (>6 weeks, often multifactorial).

\section*{Pathophysiology}
Acute paronychia typically follows disruption of the seal between the nail plate and proximal nail fold, allowing bacterial inoculation. \textit{S. aureus} is the predominant pathogen; streptococci and \textit{Pseudomonas} also occur. Chronic paronychia results from persistent irritant dermatitis with repeated wetting and maceration, disrupting the cuticle barrier. Secondary colonization by \textit{Candida albicans}, \textit{Pseudomonas}, and enteric Gram-negative rods perpetuates inflammation.

\section*{Etiology}
\begin{itemize}
  \item \textbf{Acute:} \textit{Staphylococcus aureus} (most common), \textit{Streptococcus pyogenes}, \textit{Pseudomonas aeruginosa}, \textit{Eikenella corrodens} (oral flora from nail biting/sucking), herpes simplex (herpetic whitlow)
  \item \textbf{Chronic:} Repeated hand wetting (dishwashers, bartenders, nurses, swimmers), chemical irritants, skin picking, immunosuppression, diabetes mellitus
  \item \textbf{Underlying dermatoses:} Psoriasis, eczema, contact dermatitis, lichen planus (predispose to chronic paronychia)
  \item \textbf{Iatrogenic:} Targeted therapies (EGFR inhibitors, MEK inhibitors) commonly cause paronychia as a class effect
\end{itemize}

\section*{Indications for Evaluation}
Seek care if:
\begin{itemize}
  \item Abscess formation with fluctuance, severe pain, or purulent drainage
  \item Proximal extension of erythema onto the digit or hand (cellulitis or flexor tenosynovitis concern)
  \item Fever, lymphangitis, or systemic signs of infection
  \item Paronychia persists >6 weeks despite conservative measures
  \item Recurrent acute paronychia or diabetes/immunocompromised state
\end{itemize}

\section*{Related Symptoms}
\begin{itemize}
  \item Felon (pulp space infection), herpetic whitlow, onycholysis, nail dystrophy, cellulitis
\end{itemize}
\textbf{Note:} Herpetic whitlow (HSV infection of the fingertip) mimics bacterial paronychia but presents with vesicular lesions, serous drainage, and severe lancinating pain; incision and drainage may cause viremia and is contraindicated.

\section*{Self-Care / Home Management}
\begin{itemize}
  \item Warm water soaks with Epsom salts or dilute vinegar (1:4 with water) for 15 minutes TID
  \item Keep hands dry; wear gloves for wet work; apply moisturizer to cuticles
  \item Avoid nail biting, cuticle picking, and aggressive manicuring
\end{itemize}

\section*{Diagnostic Approach}
\begin{itemize}
  \item Clinical diagnosis: Erythema, swelling, and tenderness of the proximal/lateral nail fold with or without fluctuance
  \item Assessment of fluctuance: If present, incision and drainage with culture of purulent material
  \item Chronic paronychia: KOH preparation and fungal culture; consider bacterial culture; evaluate for underlying dermatosis
  \item Tzanck smear or viral culture if vesicles present (suspected herpetic whitlow)
\end{itemize}

\section*{Treatment / Management Overview}
\begin{itemize}
  \item Acute without abscess: Topical antibiotic (mupirocin TID, gentamicin) plus corticosteroid (hydrocortisone 1\% BID)
  \item Acute with abscess: Incision and drainage with \#11 blade, blunt dissection, and packing; oral cephalexin 500 mg QID or clindamycin 300 mg TID for 5 days
  \item Chronic: Keep hands dry (gloves), topical corticosteroid (betamethasone 0.05\%) BID plus antifungal (clotrimazole) BID for 4--6 weeks
  \item Systemic: Oral antifungal (fluconazole 200 mg weekly or itraconazole pulse therapy) for recalcitrant chronic paronychia
  \item Medication-related: Emollients and topical steroids for EGFR inhibitor-induced paronychia; dose modification of causative drug
\end{itemize}

\clearpage
